\documentclass[11pt,a4paper]{article}

%========================
% PACKAGES
%========================
\usepackage[utf8]{inputenc}
\usepackage[T1]{fontenc}
\usepackage[french]{babel}
\usepackage[a4paper,margin=2cm]{geometry}
\usepackage{amsmath,amssymb}
\usepackage{lmodern}
\usepackage{fancyhdr}
\usepackage{lastpage}
\usepackage{enumitem}
\usepackage{array}
\usepackage{xcolor}
\usepackage{tcolorbox}

%========================
% MISE EN PAGE
%========================
\setlength{\headheight}{14pt}
\pagestyle{fancy}
\fancyhf{}
\fancyhead[L]{Mathématiques -- Troisième}
\fancyhead[C]{Devoir surveillé}
\fancyhead[R]{Géométrie}
\fancyfoot[C]{Page \thepage/\pageref{LastPage}}

\setlist[enumerate]{leftmargin=1cm}

\newcommand{\points}[1]{\hfill \textbf{[#1 pts]}}
\newcommand{\totalexo}[1]{\begin{flushright}\textbf{Total exercice : #1 points}\end{flushright}}

\begin{document}

%========================
% TITRE
%========================

\begin{center}
    {\Large \textbf{DEVOIR SURVEILLÉ DE MATHÉMATIQUES}}\\[0.2cm]
    {\large \textbf{Géométrie -- Niveau 3e}}\\[0.2cm]
    \textbf{Durée conseillée : 1 h 30}\\[0.2cm]
    \textbf{Total : 47 points + 6 points bonus}
\end{center}

\vspace{0.4cm}

\begin{tcolorbox}[colback=gray!5,colframe=black,title=\textbf{Consignes générales}]
\begin{itemize}
    \item La calculatrice est autorisée.
    \item Toute réponse doit être justifiée.
    \item Les figures peuvent être refaites à main levée, mais les codages doivent être respectés.
    \item Les démonstrations doivent être rédigées avec les données, le théorème utilisé et la conclusion.
    \item Les longueurs seront données en cm, sauf indication contraire.
    \item Lorsqu'un résultat est demandé sous forme approchée, on donnera une valeur arrondie au dixième.
\end{itemize}
\end{tcolorbox}

\vspace{0.5cm}

%====================================================
\section*{Exercice 1 -- Applications directes : Pythagore et trigonométrie}

\textbf{Cet exercice vérifie les bases du cours. Les questions sont indépendantes.}

\medskip

\begin{enumerate}
    \item Le triangle $ABC$ est rectangle en $A$. On donne :
    \[
    AB=6 \quad \text{et} \quad AC=8.
    \]
    \begin{enumerate}
        \item Préciser quel côté est l'hypoténuse. \points{1}
        \item Écrire l'égalité de Pythagore dans le triangle $ABC$. \points{1,5}
        \item Calculer $BC$. \points{2}
    \end{enumerate}

    \medskip

    \item Le triangle $DEF$ est rectangle en $D$. On donne :
    \[
    DE=5 \quad \text{et} \quad EF=13.
    \]
    \begin{enumerate}
        \item Préciser quel côté est l'hypoténuse. \points{1}
        \item Calculer $DF$. \points{2,5}
    \end{enumerate}

    \medskip

    \item Dans un triangle $RST$ rectangle en $R$, on considère l'angle $\widehat{RST}$. On donne :
    \[
    RS=7 \quad \text{et} \quad ST=12.
    \]
    \begin{enumerate}
        \item Pour l'angle $\widehat{RST}$, dire quel côté est adjacent et quel côté est l'hypoténuse. \points{1}
        \item Écrire une formule trigonométrique permettant de calculer $\widehat{RST}$. \points{1,5}
        \item Calculer $\widehat{RST}$, arrondi au degré près. \points{2}
    \end{enumerate}
\end{enumerate}

\totalexo{12,5}

%====================================================
\section*{Exercice 2 -- Choisir le bon théorème}

On considère la figure suivante, que l'on pourra refaire à main levée.

\medskip

Les points $A$, $B$, $M$ sont alignés dans cet ordre.  
Les points $A$, $C$, $N$ sont alignés dans cet ordre.  
Les droites $(MN)$ et $(BC)$ sont parallèles.

On donne :
\[
AM=4,8, \quad AB=12, \quad AC=15, \quad MN=5,6.
\]

\begin{enumerate}
    \item Faire une figure à main levée en respectant les alignements et le parallélisme. \points{1}

    \item Quel théorème peut-on utiliser pour calculer $AN$ ? Justifier en citant les conditions nécessaires. \points{3}

    \item Calculer $AN$. \points{3}

    \item Calculer $BC$. \points{3}

    \item On donne maintenant un autre triangle $PQR$ tel que :
    \[
    PQ=9, \quad PR=12, \quad QR=15.
    \]
    Montrer que le triangle $PQR$ est rectangle. Préciser en quel point. \points{4}
\end{enumerate}

\totalexo{14}

%====================================================
\section*{Exercice 3 -- Raisonnement géométrique rédigé}

On considère un quadrilatère $ABCD$.

Les diagonales $[AC]$ et $[BD]$ se coupent en un point $O$.

On sait que :
\[
OA=OC, \qquad OB=OD.
\]

De plus, on sait que :
\[
AB=6, \qquad BC=8, \qquad AC=10.
\]

\begin{enumerate}
    \item Que représente le point $O$ pour le segment $[AC]$ ? Justifier. \points{1,5}

    \item Que représente le point $O$ pour le segment $[BD]$ ? Justifier. \points{1,5}

    \item Démontrer que le quadrilatère $ABCD$ est un parallélogramme. \points{3}

    \item Dans le triangle $ABC$, démontrer que le triangle $ABC$ est rectangle. \points{4}

    \item En déduire que le parallélogramme $ABCD$ est un rectangle. Justifier soigneusement. \points{3}
\end{enumerate}

\totalexo{13}

%====================================================
\section*{Exercice 4 -- Problème contextualisé : une rampe d'accès}

Une mairie souhaite installer une rampe d'accès pour permettre l'entrée dans un bâtiment.

La porte d'entrée est située à une hauteur de $0,75$ m par rapport au sol.  
Pour respecter une inclinaison raisonnable, la rampe doit partir d'un point situé à $6$ m horizontalement de la porte.

On modélise la situation par un triangle rectangle $ABC$ :

\begin{itemize}
    \item $A$ est le point au sol au départ de la rampe ;
    \item $B$ est le point au sol situé verticalement sous l'entrée ;
    \item $C$ est le seuil de la porte ;
    \item le triangle $ABC$ est rectangle en $B$ ;
    \item $AB=6$ m et $BC=0,75$ m ;
    \item la rampe correspond au segment $[AC]$.
\end{itemize}

\begin{enumerate}
    \item Faire un schéma à main levée de la situation. \points{1}

    \item Calculer la longueur $AC$ de la rampe. Donner une valeur arrondie au centième de mètre. \points{3}

    \item On appelle $\alpha$ l'angle formé par la rampe avec le sol, c'est-à-dire l'angle $\widehat{BAC}$.
    \begin{enumerate}
        \item Dans le triangle rectangle $ABC$, exprimer $\tan(\alpha)$. \points{2}
        \item Calculer $\alpha$, arrondi au dixième de degré. \points{2}
    \end{enumerate}

    \item Pour être acceptée, la rampe doit former avec le sol un angle inférieur à $8^\circ$.
    
    La rampe proposée respecte-t-elle cette condition ? Justifier. \points{2}

    \item La mairie envisage finalement de construire une rampe de longueur exactement $8$ m, toujours pour atteindre une hauteur de $0,75$ m.
    
    Calculer la distance horizontale nécessaire au sol. Donner une valeur arrondie au centième de mètre. \points{4}
\end{enumerate}

\totalexo{14}

%====================================================
\section*{Exercice 5 -- Exercice difficile : retrouver une erreur et corriger un raisonnement}

Un élève veut démontrer que deux droites sont parallèles dans la configuration suivante.

Les points $A$, $B$, $M$ sont alignés dans cet ordre.  
Les points $A$, $C$, $N$ sont alignés dans cet ordre.

On donne :
\[
AM=6, \quad AB=10, \quad AN=9, \quad AC=15.
\]

L'élève écrit :

\begin{tcolorbox}[colback=gray!5,colframe=black]
On calcule :
\[
\frac{AM}{AC}=\frac{6}{15}=0,4
\]
et
\[
\frac{AN}{AB}=\frac{9}{10}=0,9.
\]
Les rapports ne sont pas égaux, donc les droites $(MN)$ et $(BC)$ ne sont pas parallèles.
\end{tcolorbox}

\begin{enumerate}
    \item Expliquer précisément l'erreur commise par l'élève. \points{3}

    \item Écrire les deux rapports corrects à comparer pour utiliser la réciproque du théorème de Thalès. \points{2}

    \item Calculer ces deux rapports. \points{2}

    \item Que peut-on conclure sur les droites $(MN)$ et $(BC)$ ? Rédiger correctement la démonstration. \points{4}

    \item On ajoute maintenant que $BC=20$.
    
    En utilisant le résultat précédent, calculer $MN$. \points{3}
\end{enumerate}

\totalexo{14}

%====================================================
\section*{Exercice 6 -- Défi bonus : une figure à analyser}

\textbf{Cet exercice est un bonus. Il demande une stratégie. Toute démarche cohérente sera valorisée.}

\medskip

Dans un repère orthonormé, on donne les points :
\[
A(0;0), \quad B(8;0), \quad C(8;6), \quad D(0;6).
\]

On place un point $M$ sur le segment $[AB]$ tel que :
\[
AM=x.
\]

On place ensuite un point $N$ sur le segment $[CD]$ tel que :
\[
DN=x.
\]

Ainsi, lorsque $x$ varie, les points $M$ et $N$ se déplacent en même temps sur deux côtés opposés du rectangle $ABCD$.

\begin{enumerate}
    \item Faire une figure à main levée. \points{1}

    \item Exprimer les coordonnées de $M$ en fonction de $x$. \points{1}

    \item Exprimer les coordonnées de $N$ en fonction de $x$. \points{2}

    \item Montrer que la longueur $MN$ ne dépend pas de $x$. \points{4}

    \item Calculer la longueur $MN$. \points{2}
\end{enumerate}

\totalexo{Bonus : 10 points}

\vspace{0.5cm}

\begin{center}
    \Large \textbf{Fin du sujet}
\end{center}

\end{document}